% Introduction: motivation, VVUQ primer, VVUQ-01 lineage, scope, contributions.
% Built by the future pass from the instructions below. Do not process now.

[Open the introduction by defining verification, validation, and uncertainty
quantification in the plain language of an oncology clinical trial, then make the
single-humanoid argument, then trace the lineage from the prior paper, then state
the scope and contributions. Aim for several connected paragraphs of publication
quality; do not enumerate file names in the prose, push those into the Methods
tables, and cite rather than describe the literature.]

[Paragraph 1, the VVUQ primer. Frame the three questions for a reader who knows
oncology trials but not credibility engineering: verification asks did we build
the code correctly, validation asks did we build the right model, and uncertainty
quantification asks how confident we are in the results. Ground this framing in
\texttt{papers/VVUQ-02/codegen/docs/vvuq\_methodology.md} and the survey vocabulary
of \texttt{sel2025vvuq}, and note the oncology precedent for VVUQ on
treatment-planning models with \texttt{zaid2025pbpk}. Connect the three questions
to the gate primitives in \texttt{papers/VVUQ-02/codegen/src/vvuq/} (verification.py,
validation.py, uncertainty.py) so the reader sees that each question is a concrete,
executed check rather than an abstraction.]

[Paragraph 2, the thesis and why one humanoid raises the bar. State the thesis
verbatim in spirit from \texttt{papers/VVUQ-02/codegen/README.md} and
\texttt{papers/VVUQ-02/execution/README.md}: because one autonomous humanoid
concentrates all error potential into two hands, with no second cart and no human
bedside assistant in the loop, the assurance layer must be larger, stricter, and
more thoroughly exercised than the control code, and any candidate behavior must
clear ten distinct humanoid-specific gates before it may ship. Make explicit that
code generation is fast and deterministic while the decision to ship is the
substantial work, and that this asymmetry is the paper's central claim.]

[Paragraph 3, the lineage from VVUQ-01. Summarise, without copying text, the
developments in \texttt{papers/VVUQ-01/final-paper/} that led here: the prior work
established a two-stage VVUQ verification-over-generation framework for an LLM
oncology code workflow (cite \texttt{kawchak2026vvuq01}). State what is new in this
paper: the framework is specialised from generic code verification to ten
humanoid-specific surgical gates, bound to a wired corpus of robotic-surgery and
service-robot safety standards, and exercised over a concrete 60-second Whipple
procedure. Reference the broader autonomous physical AI oncology program with
\texttt{kawchak2026sponsor}, \texttt{kawchak2026federated}, and
\texttt{kawchak2026prediction}, and the PDAC modelling lineage with
\texttt{kawchak2025qsppdac} and \texttt{kawchak2025pdactwin}.]

[Paragraph 4, scope and the regulatory moment. State the frozen scope from
\texttt{papers/VVUQ-02/codegen/README.md} (one simulated patient PAT-PDAC-0001, one
classic pancreaticoduodenectomy with pylorus preservation, one robot, the
hypothetical 2030 Unitree H2-Surgical 1.0 with two seven-DOF arms and two
twenty-DOF dexterous hands, and a 32-iteration deterministic design at seed
20260525). Connect to the regulatory moment respectfully but opportunistically:
the FDA real-time clinical trials initiative (\texttt{fda2026realtime}) and the FDA
computational-model credibility guidance (\texttt{fda2023credibility}) make a fast,
standards-anchored, autonomous assurance pipeline timely, and argue that the United
States should lead the world in surgical-humanoid VVUQ tied to external standards
for patient safety, efficacy, and trial speed. Keep all references to the FDA and
other bodies respectful and non-presumptuous.]

[Paragraph 5, contributions and roadmap. List three to four contributions mapped
to the Results section, drawn from the \texttt{Execution results summary} and
\texttt{Key quantitative results to cite} tables in
\texttt{papers/VVUQ-02/execution/README.md}: the ten-gate decision surface, the
32-iteration sweep and four-entrant tournament, the featured 1000-row sensor stream
and 60-second timeline, and the reproducible standards-anchored credibility
argument. End with a one-sentence roadmap of the paper. Do not include a table in
this section; the Introduction is prose only, and the contributions table belongs
to Methods.]
